\section{Fundamentação teórica e trabalhos relacionados}\label{sec:fundamentacao}
\subsection{Modelos de linguagem}
Os LLMs atuais derivam, em grande parte, da arquitetura Transformer, baseada em mecanismos de
atenção~\citep{vaswani2017attention}. Esses modelos são pré-treinados em grandes volumes de texto e
modelam diretamente a probabilidade de sequências de texto~\citep{liu2021pretrain}. Assim, podem
ser usados em novas tarefas sem alterar os parâmetros, apenas mudando a entrada~\citep{sahoo2024survey}.
É essa possibilidade que separa o foco deste trabalho das abordagens baseadas em treinamento ou ajuste fino.

\subsection{Engenharia de prompt}
Um prompt é a instrução dada ao modelo para impor regras, automatizar processos e obter certas
qualidades na saída~\citep{white2023patterns}. Engenharia de prompt é a prática de projetar essas
entradas para obter o comportamento desejado de um modelo já treinado~\citep{sahoo2024survey}.
\citetexto{White et al.}{white2023patterns} tratam prompts como uma forma de programação e catalogam
padrões reutilizáveis, por analogia com os padrões de software.

\subsection{Técnicas de prompting}
As técnicas abaixo ilustram o tipo de objeto que a revisão pretende classificar. A lista não é
exaustiva, sua apresentação aqui não indica qual funciona melhor, e a menção não decide a
elegibilidade de nenhuma família de estudos.
\begin{itemize}
  \item \textit{Zero-shot}: apenas uma instrução em linguagem natural, sem
  demonstrações~\citep{brown2020fewshot}.
  \item \textit{Few-shot}: o prompt inclui algumas demonstrações da tarefa.
  \citetexto{Brown et al.}{brown2020fewshot} avaliam essa condição, ao lado de \textit{one-shot} e \textit{zero-shot},
  como formas de medir o aprendizado em contexto, de modo que os dois termos não são sinônimos perfeitos.
  \item \textit{Chain-of-thought}: na formulação de \citetexto{Wei et al.}{wei2022cot}, as demonstrações
  trazem uma série de passos intermediários de raciocínio. Variantes sem demonstrações só serão
  descritas se aparecerem no corpus.
  \item ReAct: o modelo intercala traços de raciocínio e ações, o que permite consultar fontes ou
  ambientes externos~\citep{yao2023react}.
  \item \textit{Self-planning}: o modelo primeiro planeja a solução e depois escreve o código,
  a partir do plano~\citep{jiang2025selfplanningv5}.
  \item \textit{Self-debugging}: demonstrações ensinam o modelo a depurar o próprio programa,
  examinando resultados de execução e explicando o código~\citep{chen2023selfdebug}.
\end{itemize}
Esses itens não são do mesmo tipo: alguns descrevem uma forma de montar o prompt, e outros, uma
sequência de interação com o modelo.
Por isso, a síntese usa duas unidades. A técnica, com o nome dado pelos autores, é a unidade das comparações.
O mecanismo é a unidade da classificação: cada técnica recebe um ou mais rótulos de um dicionário fechado de onze
mecanismos, que são instrução, contexto, restrições, exemplos, decomposição, raciocínio, \textit{feedback},
amostragem, multiagente, otimização de prompt e outro. Assim, uma sequência de interação como o
\textit{self-debugging} aparece como combinação de mecanismos, por exemplo \textit{feedback}, raciocínio e
exemplos, e não como categoria própria. A classificação é multirrótulo e distingue os mecanismos do pacote
completo dos que a técnica acrescenta ao prompt base.

\subsection{Tarefas e formas de avaliação}
Neste trabalho, geração é a produção de código a partir de uma descrição, especificação ou exemplo.
Tradução é a conversão de código entre linguagens ou plataformas, incluindo migração e transpilação.
Refatoração é a melhoria da estrutura do código sem mudança de comportamento. Depuração e reparo
consistem em localizar e corrigir defeitos. Testes entram quando apoiam essas tarefas. Essas
definições delimitam o escopo da revisão, mas não resolvem todos os casos de fronteira: a
Seção~\ref{sec:metodologia} mostra a calibração feita e os casos ainda em aberto.

A avaliação costuma recorrer a \textit{benchmarks}. O HumanEval mede a correção funcional de
programas sintetizados a partir de \textit{docstrings}~\citep{chen2021codex}. O SWE-bench usa
problemas reais do GitHub, em que o modelo precisa editar uma base de código para resolver uma
\textit{issue}~\citep{jimenez2024swebench}. Resultados nesses conjuntos não demonstram, por si,
ganhos de produtividade, segurança ou custo. \citetexto{Sallou et al.}{sallou2024breaking} discutem
ameaças à validade no uso de LLMs em engenharia de software, e parte delas orienta a avaliação de
qualidade prevista neste trabalho.

\subsection{Trabalhos relacionados}
% AFIRMACAO: A119
\citetexto{Hou et al.}{hou2024slr} revisaram 395 estudos sobre LLMs em engenharia de software, de janeiro de 2017 a
janeiro de 2024. No prompting, descrevem as técnicas e contam sua frequência: \textit{few-shot} em 88 estudos,
\textit{zero-shot} em 79 e CoT em 18.
% AFIRMACAO: A120
\citetexto{Fan et al.}{fan2023survey} fazem um levantamento da área, e não uma revisão sistemática, e apontam as técnicas
híbridas, que combinam engenharia de software tradicional e LLMs, como centrais para resultados confiáveis.
Ambos abrangem muitas atividades de engenharia de software, e não apenas prompting.
% AFIRMACAO: A121
Em sentido oposto, \citetexto{Sahoo et al.}{sahoo2024survey} organizam 41 técnicas de prompting por área de aplicação,
entre elas técnicas próprias de código, como \textit{Program of Thoughts}, SCoT e \textit{Chain of Code}, e
\citetexto{Liu et al.}{liu2021pretrain} sistematizam o paradigma de prompting em processamento de linguagem natural, sem se
limitar a tarefas de código.

A busca deste trabalho também encontrou estudos secundários próximos, registrados como apoio.
% AFIRMACAO: A122
\citetexto{Yang et al.}{yang2025aprsurvey} (S820) classificam 66 sistemas de reparo baseados em LLM em ajuste fino,
prompting, procedimentos e agentes e registram, para o resultado principal de cada sistema, o \textit{benchmark}, a
métrica, o modelo e a suposição sobre a localização do defeito.
% AFIRMACAO: A123
\citetexto{Masthan Ali et al.}{masthanali2026energy} (S1246) reúnem 18 estudos sobre a eficiência energética do código
gerado por LLMs e concluem que nenhuma estratégia de prompting a melhora de forma consistente entre configurações.

Este trabalho cruza os dois recortes: técnicas de prompting nas quatro tarefas, com extração do tipo de
contrafactual, da direção de cada comparação, das condições de avaliação, dos riscos e das mitigações. Os textos
completos dessas revisões foram lidos, nas versões do arXiv indicadas nas referências, para a comparação da Seção~\ref{sec:comparacao}. Como a
busca não foi desenhada para localizar revisões, não se afirma que essa combinação seja inédita.
